%%%%%%%%%%%%%%%%%%%%%%%%
% $Autor: MansourAhmadyPhoulady, RanjeetsingTawar, HemanthPrabhakaran$
% $Pfad: Car_Licence_Plate_Identification_with_a_JetsonNano$
% $Version: 4.0.0 $
%
% !TeX encoding = utf8
%
%%%%%%%%%%%%%%%%%%%%%%%%


\chapter{Software Tests }
\label{chapter:SoftwareTests}

	The documentation for the car license plate detection project, designed for the Jetson Nano, offers thorough descriptions of each module, function, and class. It illuminates the objectives, capabilities, and significance of various components within the project. The documentation further includes test classes and methods to verify the accuracy and dependability of the features implemented. Organizing the project by modules aids in maintaining code modularity and ease of upkeep.
	
	This detailed documentation acts as an essential guide for understanding the project, facilitating teamwork, and guiding future enhancements. It increases transparency, enabling developers and stakeholders to effectively comprehend the project's structure and operational mechanisms.


\section{Functions}

\subsection{\texttt{\_\_init\_\_(self, pipeline\_config\_path, checkpoint\_path, label\_map\_path)}}

\subsubsection{Description}
The \texttt{\_\_init\_\_} method acts as the constructor for the \texttt{WebcamObjectDetection} class. It initializes the object detection model with necessary configurations, restores the model from a specified checkpoint, and prepares it for detection tasks.

\subsubsection{Parameters}
\begin{itemize}
	\item \texttt{pipeline\_config\_path}: Path to the pipeline configuration file, detailing model architecture and input specifications.
	\item \texttt{checkpoint\_path}: Path to the saved model checkpoint, containing the pre-trained model weights.
	\item \texttt{label\_map\_path}: Path to the label map file, associating each category with an integer ID.
\end{itemize}

\subsubsection{Key Operations}
\begin{itemize}
	\item Model building and loading from the provided configuration.
	\item Checkpoint restoration for weight initialization.
	\item Label map loading for human-readable label mapping.
\end{itemize}

\subsection{\texttt{detect\_fn(self, image)}}

\subsubsection{Description}
The \texttt{detect\_fn} method, optimized with \texttt{tf.function}, performs object detection on a single image. It preprocesses the image, runs it through the detection model, and post-processes the output to return detection results.

\subsubsection{Parameters}
\begin{itemize}
	\item \texttt{image}: A TensorFlow tensor of the image on which detection is to be performed.
\end{itemize}

\subsubsection{Key Operations}
\begin{itemize}
	\item Image preprocessing as required by the model.
	\item Prediction execution through the detection model.
	\item Post-processing of predictions to yield detection results.
\end{itemize}

\subsection{\texttt{run(self)}}

\subsubsection{Description}
The \texttt{run} method orchestrates capturing video from the webcam, performing object detection on each frame, and visualizing the results in real-time, until the process is terminated by the user.

\subsubsection{Key Operations}
\begin{itemize}
	\item Initialization of video capture and frame capturing from the webcam.
	\item Execution of detection on each frame and processing of the detection results.
	\item Visualization of detection results on the frames.
	\item Display of annotated frames and management of user termination requests.
\end{itemize}

\section{Classes}

\subsection{Detailed Description}

The \textbf{WebcamObjectDetection} class is engineered to enable real-time object detection using a webcam feed as input. It utilizes TensorFlow's Object Detection API to preprocess video frames, execute detections, and annotate the captured frames with detection results in real time. This class simplifies the complexity involved in interfacing with TensorFlow's more granular APIs by offering an easy-to-use interface for conducting object detection tasks.

\subsection{Core Functionalities}

\begin{itemize}
	\item \textbf{Initialization and Configuration:} The class initializes the TensorFlow detection model upon creation, loading necessary configurations from a specified pipeline configuration file, restoring pre-trained weights from a checkpoint, and arranging the label map for converting detection class IDs into human-readable labels.
	
	\item \textbf{Detection Function:} It defines an optimized TensorFlow function, \texttt{detect\_fn}, responsible for processing an input image tensor through the detection model, including pre and post-processing steps, to return bounding boxes, class IDs, and detection scores.
	
	\item \textbf{Real-time Webcam Stream Processing:} Through the \texttt{run} method, the class captures video frames from a webcam in real-time, applies object detection on each frame, and visualizes the results on the live video, creating a continuous feedback loop of detection outputs until the process is manually terminated.
\end{itemize}

\subsection{Key Methods}

\begin{description}
	\item[\texttt{\_\_init\_\_(self, pipeline\_config\_path, checkpoint\_path, label\_map\_path)}:] Initializes the detection environment with configurations, model weights, and label mapping.
	
	\item[\texttt{detect\_fn(self, image)}:] Executes the detection process on a given image tensor, encapsulating image preprocessing, model prediction, and prediction post-processing.
	
	\item[\texttt{run(self)}:] Starts the webcam feed, processes video frames in real-time for object detection, and visualizes detection outcomes on the streamed video frames.
\end{description}

\subsection{Usage Scenario}

The \textbf{WebcamObjectDetection} class is particularly useful for developers and researchers requiring an efficient method to integrate real-time object detection into applications such as surveillance systems, interactive installations, or any context necessitating live video stream analysis. It abstracts away the complexities of TensorFlow's Object Detection API, providing a straightforward path to embedding object detection functionalities.

\subsection{Technical Details}

Leveraging TensorFlow's optimized functions and best practices for inference, the class is designed to ensure the detection pipeline is sufficiently efficient for processing video frames with minimal latency, suitable for real-time applications. Direct visualization of detection results on the video stream offers instant insight into model performance, facilitating rapid adjustments or tuning as needed.

	
\section{Test File Descriptions}

\subsection{File 1: \texttt{test\_initialization.py}}
This test file verifies the initialization process of the \texttt{WebcamObjectDetection} class, ensuring that the model and configurations are properly loaded.

\begin{lstlisting}[language=Python]
	from webcam_object_detection import WebcamObjectDetection
	
	def test_initialization():
	pipeline_config_path = 'path/to/your/pipeline.config'
	checkpoint_path = 'path/to/your/checkpoint/ckpt-0'
	label_map_path = 'path/to/your/label_map.pbtxt'
	
	try:
	detector = WebcamObjectDetection(pipeline_config_path, checkpoint_path, label_map_path)
	print("Initialization Successful")
	except Exception as e:
	print(f"Initialization Failed: {e}")
	
	if __name__ == "__main__":
	test_initialization()
\end{lstlisting}

\subsection{File 2: \texttt{test\_detect\_function.py}}
This script tests the \texttt{detect\_fn} method, confirming that it can successfully process an image and return detection results.

\begin{lstlisting}[language=Python]
	import cv2
	import tensorflow as tf
	from webcam_object_detection import WebcamObjectDetection
	
	def test_detect_function():
	pipeline_config_path = 'path/to/your/pipeline.config'
	checkpoint_path = 'path/to/your/checkpoint/ckpt-0'
	label_map_path = 'path/to/your/label_map.pbtxt'
	
	detector = WebcamObjectDetection(pipeline_config_path, checkpoint_path, label_map_path)
	
	# Load an image
	image_path = 'path/to/test/image.jpg'
	image_np = cv2.imread(image_path)
	input_tensor = tf.convert_to_tensor([image_np], dtype=tf.float32)
	
	detections = detector.detect_fn(input_tensor)
	print("Detection Results:", detections)
	
	if __name__ == "__main__":
	test_detect_function()
\end{lstlisting}

\subsection{File 3: \texttt{test\_run\_method.py}}
This file tests the \texttt{run} method, assessing its ability to capture video from a webcam, detect objects in real-time, and display the annotated video stream.

\begin{lstlisting}[language=Python]
	from webcam_object_detection import WebcamObjectDetection
	
	def test_run_method():
	pipeline_config_path = 'path/to/your/pipeline.config'
	checkpoint_path = 'path/to/your/checkpoint/ckpt-0'
	label_map_path = 'path/to/your/label_map.pbtxt'
	
	detector = WebcamObjectDetection(pipeline_config_path, checkpoint_path, label_map_path)
	detector.run()
	
	if __name__ == "__main__":
	test_run_method()
\end{lstlisting}

\section{Conclusion}
The provided test files offer a basic framework for validating the functionalities of the \texttt{WebcamObjectDetection} class. They are crucial for ensuring the reliability and effectiveness of the object detection process in real-time applications.

		
		\section{ Automation}
		
\section{Automation 1: Model Training and Evaluation Automation}

\subsection{Python Pseudo-Code}

\begin{lstlisting}[language=Python]
	import subprocess
	import tensorflow as tf
	from object_detection import model_main_tf2
	
	def train_and_evaluate_model(pipeline_config_path, model_dir, num_train_steps):
	"""
	Automates the training and evaluation process of an object detection model.
	
	Args:
	pipeline_config_path (str): Path to the pipeline configuration file.
	model_dir (str): Directory where model checkpoints will be saved.
	num_train_steps (int): Number of training steps.
	"""
	command = f"python model_main_tf2.py --model_dir={model_dir} --pipeline_config_path={pipeline_config_path} --num_train_steps={num_train_steps} --alsologtostderr"
	subprocess.run(command.split())
	
	# Example usage
	train_and_evaluate_model('path/to/your/pipeline.config', 'path/to/your/model_dir', 10000)
\end{lstlisting}

\subsection{Explanation}
This script automates the process of training and evaluating an object detection model using TensorFlow's Object Detection API. The \texttt{train\_and\_evaluate\_model} function takes the path to the pipeline configuration, the directory for saving model checkpoints, and the number of training steps as arguments. It then constructs and executes a command to start the training and evaluation process. This automation simplifies the execution of training sessions, especially when experimenting with different configurations or datasets.

\section{Automation 2: Dataset Preparation and Annotation}

\subsection{Python Pseudo-Code}

\begin{lstlisting}[language=Python]
	# Assuming the use of an existing tool or script for dataset conversion
	
	def prepare_dataset(images_dir, annotations_dir, output_dir):
	"""
	Automates the conversion of datasets to TFRecord format.
	
	Args:
	images_dir (str): Directory containing images.
	annotations_dir (str): Directory containing annotations.
	output_dir (str): Output directory for TFRecords.
	"""
	command = f"python dataset_converter.py --images_dir={images_dir} --annotations_dir={annotations_dir} --output_dir={output_dir}"
	subprocess.run(command.split())
	
	# Example usage
	prepare_dataset('path/to/images', 'path/to/annotations', 'path/to/output')
\end{lstlisting}

\subsection{Explanation}
This pseudo-code demonstrates how to automate the dataset preparation process, assuming the existence of a script (\texttt{dataset\_converter.py}) that converts datasets into the TFRecord format required by TensorFlow. The \texttt{prepare\_dataset} function simplifies the conversion process, enabling quick preparation of datasets for training.


\chapter{Hardware Tests}


\section{Hardware Parts}
\begin{itemize}[label={$\bullet$}]
	\item \textbf{Power Supply:} 5V 4A DC power supply with a 5.5mm barrel connector.
	\item \textbf{SD Card:} MicroSD card with at least 16GB capacity, flashed with the Jetson Nano's operating system.
	\item \textbf{USB Keyboard and Mouse:} Required for input during setup.
	\item \textbf{USB Camera:} Compatible USB camera for testing purposes.
	\item \textbf{Keyboard and Mouse:} Connect USB keyboard and mouse to the Jetson Nano's USB ports.
\end{itemize}

\subsection*{Initialization Steps}
\begin{itemize}[label={$\bullet$}]
	\item \textbf{Language:} Follow on-screen prompts to select preferred language during setup.
	\item \textbf{Time:} Set system time during initial setup process.
\end{itemize}

\subsection*{Setting up Jetson Nano}
\begin{enumerate}[label={\arabic*.}]
	\item Insert microSD card with pre-flashed OS into Nano's card slot.
	\item Connect power supply to Nano.
	\item Connect monitor to Nano's HDMI port.
	\item Power on Jetson Nano by plugging in power supply.
\end{enumerate}

\subsection*{Image Testing}
Once Jetson Nano boots up:
\begin{itemize}[label={$\bullet$}]
	\item Test functionality using various applications or scripts.
	\item Utilize USB camera for image capture and processing with Nano's GPU capabilities.
	\item Install additional software or libraries as needed for testing requirements.
\end{itemize}

\textbf{Note:} Follow instructions provided by NVIDIA for setup and configuration. Ensure access to necessary documentation for troubleshooting.

\section{Hardware Functions}

\begin{enumerate}[label=\textbf{\arabic*.}]
	\item \textbf{Power Supply:}
	\begin{itemize}
		\item The power supply provides electrical power to the Jetson Nano board. It typically connects via a micro-USB or DC barrel jack connector.
		\item \textbf{Function:} Powers the board, allowing it to function properly.
	\end{itemize}
	
	\item \textbf{SD Card:}
	\begin{itemize}
		\item The SD card is used to store the operating system and other software required for the Jetson Nano to operate.
		\item \textbf{Function:} Provides the necessary storage for the operating system and files.
	\end{itemize}
	
	\item \textbf{Initialization Steps (Keyboard, Mouse, Language, Time):}
	\begin{itemize}
		\item When setting up the Jetson Nano for the first time, you'll typically connect a keyboard and mouse to navigate through the setup process.
		\item Language and time settings are usually configured during the initial setup process.
		\item \textbf{Function:} Enables the user to interact with the Jetson Nano and set basic configurations such as language and time.
	\end{itemize}
	
	\item \textbf{Attaching USB Camera:}
	\begin{itemize}
		\item You can attach a USB camera to the Jetson Nano to capture images or video feed for various applications like computer vision or video processing.
		\item \textbf{Function:} Provides visual input to the Jetson Nano for processing.
	\end{itemize}
	
	\item \textbf{Setting up Jetson Nano:}
	\begin{itemize}
		\item This involves installing the operating system, configuring network settings, updating software packages, and installing any additional software required for your specific use case.
		\item \textbf{Function:} Prepares the Jetson Nano for use by configuring it with the necessary software and settings.
	\end{itemize}
	
	\item \textbf{Image Testing:}
	\begin{itemize}
		\item Once the Jetson Nano is set up, you can perform image testing to verify that the system is functioning correctly. This might involve capturing images from a connected camera, processing them using software on the Jetson Nano, and verifying the output.
		\item \textbf{Function:} Validates that the Jetson Nano is correctly configured and capable of performing its intended tasks, such as image processing or computer vision.
	\end{itemize}
\end{enumerate}

\newpage

\section{Documentation}


\begin{longtable}{|p{3cm}|p{2.5cm}|p{3cm}|p{1.5cm}|p{3cm}|}
	\hline
	\textbf{Test Description} & \textbf{Expected} & \textbf{Actual} & \textbf{Test Result} & \textbf{Remarks} \\
	\hline
	\endfirsthead % Header for the first page
	
	\multicolumn{5}{c}%
	{{\bfseries \tablename\ \thetable{} -- continued from previous page}} \\
	\hline
	\textbf{Test Description} & \textbf{Expected} & \textbf{Actual} & \textbf{Test Result} & \textbf{Remarks} \\
	\hline
	\endhead % Header for the rest of the pages
	
	\hline
	\multicolumn{5}{|r|}{{Continued on next page}} \\ \hline
	\endfoot % Footer for all but the last page
	
	\endlastfoot % Footer for the last page
	
	Format the 32 GB SD card using an SD card formatter & SD card formatted & First time failed, second time success & Pass & SD Memory Card Formatted \\
	\hline
	Retrieve the image from the provided Nvidia link & Downloaded successfully & Successfully downloaded & Pass & Image retrieved \\
	\hline
	Burn or etch the image onto the SD card & Successful Flashing JetPack OS & Flashed JetPack OS on SD card & Pass & Balena Etcher Software Used \\
	\hline
	Open the etching software, Balena Etcher and select the downloaded image & Successfully loading of OS image & Successful loading of OS image & Pass & Balena Etcher Software Used \\
	\hline
	Unmount the card when it shows 100\% complete & Unmounted without corrupting SD card & Unmounted without corrupting SD card & Pass & Done Successfully \\
	\hline
	Insert SD card in Jetson Nano board & Mount SD card in Jetson Nano & Mount SD card in Jetson Nano & Pass & Done \\
	\hline
	Follow the installation steps and select username, language, keyboard, and time settings & Login into Jetson Nano developer kit & Successful login into Jetson Nano & Pass & Choose English and German time \\
	\hline
	Install the media device packages using v4l-utils & Install package dependency & Packages not installed & Failed & No internet connection, used LAN cable \\
	\hline
	Check the attached USB cameras, supported formats, and resolutions & Camera setup and capturing of image & Image captured and successful setup of camera & Pass & Typed cheese in the terminal to check camera and Lsusb to check the input details \\
	\hline
	Deploy the code in the Jetson Nano along with the environment & Code deployed successfully & Environment and code setup done accordingly & Pass &  \\
	\hline
	Running of the python file to detect license plate using Logitech webcam & Successful detection of the license plate & Successfully detected the license plate & Pass &  \\
	\hline
	Detection example on the image & Rectangle plotted in the images & Plotted green rectangle using CV2 around the license plate & Pass &  \\
	\hline
	Result of Easy-OCR on segmented image & Recognize the characters of LP & EasyOCR results generated & Pass &  \\
	\hline
	Detection Live feed example with a USB camera & Detection of plate and storing it with images & Detected the license plate and stored the text in detection\_results csv and images in Detected\_images folder & Pass &  \\
	\hline
\end{longtable}

	
		
		
		
		
		
		
		
		
		
		
		
		
		
		
		
